\section{Equidade e a Ética do Acesso}
\label{sec:equidade}

\subsection{O Risco do \textit{Digital Divide} na Odontologia}

A dimensão bioética da justiça distributiva impõe um questionamento severo sobre a incorporação de gêmeos digitais no Brasil. O acesso a essas plataformas requer conectividade de alta velocidade, processamento intensivo e hardware dispendioso (scanners IOS, tomógrafos CBCT). Existe um risco iminente de que a tecnologia, em vez de democratizar a excelência clínica, atue como um vetor de amplificação das iniquidades em saúde bucal, consolidando o que Cuadros et al. denominam "desertos digitais" \cite{cuadros2023assessing}.

Enquanto as clínicas em centros metropolitanos do eixo Sul-Sudeste adotarão gêmeos digitais preditivos para minimizar falhas implantodônticas, o sistema público nas regiões Norte e Nordeste, frequentemente desprovido de radiografia periapical digital, aprofundará seu atraso. A ética do acesso exige que as políticas públicas sejam desenhadas de modo a subsidiar a adoção tecnológica não como luxo, mas como infraestrutura básica de saúde.

\subsection{Inclusão Digital como Política de Estado}

Para que a equidade seja resguardada, a implementação tecnológica no Sistema Único de Saúde (SUS) deve ancorar-se em plataformas de código aberto. Ecossistemas abertos e descentralizados, como o OpenTwins \cite{robles2023opentwins} ou as extensões do OpenCode, reduzem os custos de licenciamento e previnem a submissão do Estado aos cartéis de software (\textit{vendor lock-in}). A tele-odontologia, amparada por gêmeos digitais, tem o potencial formidável de transportar o planejamento de alta complexidade para as comunidades isoladas \cite{frota2025sus}, contanto que haja um planejamento central voltado à conectividade (como o 5G/6G em unidades rurais).
